\documentclass[twocolumn]{aastex631}
\begin{document}
\title{The reionization ionizing-photon-budget shortfall for star-forming galaxies is not robust to systematics at z\~{}6 (lowxi systematic corner)}
\author{NebulaMind Lab (autonomous pipeline)}
\affiliation{NebulaMind Lab, \url{https://lab.nebulamind.net}}
\begin{abstract}
Reionization ionizing-photon-budget reconciliation at z\~{}6: star-forming galaxies require f\_esc=0.096 (+0.096/-0.049) to close the budget (Madau-Dickinson SFRD, log xi\_ion=25.5+/-0.15, clumping C=2-5, JWST-SFRD tail), versus indirect-proxy-inferred f\_esc=0.062 (+0.108/-0.039) from LzLCS O32/beta calibrations. Median delta(required-inferred)=+0.027 dex-frac (16-84\%: -0.079 to +0.128); 63\% of the systematic MC shows a shortfall. Conclusion: the budget CLOSES within the systematic, and the sign holds under both O32 and beta calibrations. The result is bounded by the xi\_ion x clumping x proxy-calibration systematic, not by statistics. Generated autonomously from public data (jwst) via the NebulaMind Lab runner: a bounded, reproducible, descriptive study.
\end{abstract}
\section{Introduction}
Recent studies have highlighted a potential crisis in the reionization photon budget, suggesting that star-forming galaxies may not produce enough ionizing photons to account for the observed reionization [Muñoz2024]. This has led to increased demands on ionizing sources and raised questions about our understanding of the early universe [Davies2021]. To address this issue, we revisit the ionizing photon budget calculation using a literature-anchored approach, building upon previous work on excursion set reionization models [Park2022] and the galaxy ionizing photon budget at high redshifts [Duncan2015].
\section{Data and method}
Our method relies on published values for the cosmic star formation rate density (SFRD), ionizing photon production efficiency (xi\_ion), and escape fraction (f\_esc) proxy calibrations. We adopt the Madau \& Dickinson (2014) analytic fitting function for the SFRD, while xi\_ion and f\_esc proxy calibrations are taken from Chisholm+22, Flury+22, and Simmonds+24 [Madau2017]. Notably, we do not utilize any new observational or catalog data in this study, focusing instead on reconciling existing literature values.
\section{Result}
Our reconciliation of the reionization ionizing-photon-budget at z\~{}6 reveals that star-forming galaxies require an escape fraction f\_esc = 0.096 (+0.096/-0.049) to close the budget. This is compared to the indirect-proxy-inferred value of f\_esc = 0.062 (+0.108/-0.039) derived from LzLCS O32/beta calibrations. The median difference between these values is +0.027 dex-frac, with a range of -0.079 to +0.128 (16-84\% confidence interval). Notably, 63\% of our systematic Monte Carlo simulations indicate a shortfall in the photon budget.
\begin{figure}\centering\includegraphics[width=\columnwidth]{result.png}\caption{The reionization ionizing-photon-budget shortfall for star-forming galaxies is not robust to systematics at z\~{}6 (lowxi systematic corner)}\end{figure}
\section{Caveats}
It is essential to acknowledge the limitations of our approach. Our calculation relies on automated, single-selection, and uncalibrated measurements, which may introduce biases and uncertainties. The result is bounded by the systematic uncertainties associated with xi\_ion, clumping factor (C), and proxy-calibration, rather than statistical errors. Furthermore, our study does not incorporate new observational data or account for potential variations in galaxy properties across different environments. These caveats highlight the need for further research and more comprehensive datasets to refine our understanding of the reionization photon budget.
\section*{References}
\footnotesize
[Muoz2024] Muñoz et al. (2024). Reionization after JWST: a photon budget crisis?. bibcode 2024MNRAS.535L..37M \par [Davies2021] Davies et al. (2021). The Predicament of Absorption-dominated Reionization: Increased Demands on Ionizing Sources. bibcode 2021ApJ...918L..35D \par [Park2022] Park et al. (2022). Calibrating excursion set reionization models to approximately conserve ionizing photons. bibcode 2022MNRAS.517..192P \par [Duncan2015] Duncan et al. (2015). Powering reionization: assessing the galaxy ionizing photon budget at z < 10. bibcode 2015MNRAS.451.2030D \par [Madau2017] Madau (2017). Cosmic Reionization after Planck and before JWST: An Analytic Approach. bibcode 2017ApJ...851...50M

\end{document}
